\subsection[Secondary Part 3 DIBL extraction]{Secondary Part 3 extraction: ELR/ESR sensitivity}
A second extraction was tested on the \emph{original} ten
$800$--$1000\,^\circ$C transfer sweeps to assess whether
drain-bias sensitivity can be compared when the predeclared
$10^{-7}$ constant-current threshold is not crossed.
The procedure is the same model-dependent ELR/ESR audit
described in Part~4. This additional analysis does
\emph{not} replace the constant-current results in
the main Part~3 metric table.

\paragraph{Extraction and provenance.}
The archived pairs use $V_{DS}=0.1/1.2$~V:
nodes $n3/n9$ (800$^\circ$C), $n13/n14$ (850$^\circ$C),
$n17/n18$ (900$^\circ$C), $n21/n22$ (950$^\circ$C),
and $n25/n26$ (1000$^\circ$C).
Each pair was verified to use the same SProcess
structure. A local, read-only Python extraction
parsed the original SDevice \texttt{.plt} curves.
For the low-drain bias, a five-point maximum-slope fit
to $|I_D|=aV_{GS}+b$ gives
$V_{\mathrm{TH,ELR}}=-b/a-0.05$~V.
For the high-drain bias, a separate maximum-slope fit
to $\sqrt{|I_D|}=cV_{GS}+d$ gives
$V_{\mathrm{TH,ESR}}=-d/c$. The reported shift is
\[
D_{\mathrm{ELR/ESR,app}}
=1000\,\frac{V_{\mathrm{TH,ELR}}(0.1\mathrm{\,V})-
V_{\mathrm{TH,ESR}}(1.2\mathrm{\,V})}{1.1\mathrm{\,V}}
\quad(\mathrm{mV/V}).
\]
The extraction was repeated with 3-, 5-, 7- and
9-point windows. These are \emph{different} strong-inversion
transport approximations used at the two drain biases:
a local straight-line fit does not demonstrate a
physically unique inversion threshold. Standard
threshold-extraction reviews recognize both techniques
and their potential differences
(A.~Ortiz-Conde \emph{et al.},
\emph{Microelectronics Reliability} \textbf{42},
583--596, 2002;
\href{https://doi.org/10.1016/S0026-2714(02)00027-6}{doi:10.1016/S0026-2714(02)00027-6}).

\begin{table}[H]\centering\small
\caption{Alternative apparent Part 3 DIBL from ELR/ESR.
Values are transcribed at the precision printed by the
successful NUS server extraction on 27 October 2026.
Neither fitted intercept is asserted to be a
constant-current threshold measurement.}
\begin{tabular}{@{}rrrrr@{}}\toprule
Anneal & $V_{\mathrm{TH,ELR}}$ &
$V_{\mathrm{TH,ESR}}$ & $D_{\mathrm{app}}$ &
Window range\\
($^\circ$C) & (V) & (V) & (mV/V) & (mV/V)\\\midrule
800  & 0.1913 & $-0.1929$ & 349.2 & 346.7--349.9\\
850  & 0.1876 & $-0.2784$ & 423.7 & 418.1--424.5\\
900  & 0.1937 & $-0.2997$ & 448.5 & 446.0--449.1\\
950  & 0.1961 & $-0.3477$ & 494.4 & 494.3--494.5\\
1000 & 0.1280 & $-0.9417$ & 972.5 & 970.5--973.1\\
\bottomrule\end{tabular}
\end{table}

\begin{figure}[H]\centering
\begin{tikzpicture}
\begin{axis}[
width=.91\linewidth,height=6.6cm,
xlabel={Activation anneal ($^\circ$C)},
ylabel={Drain-bias-dependent gate shift (mV/V)},
xmin=780,xmax=1020,ymin=0,ymax=1100,
xtick={800,850,900,950,1000},
grid=both,
legend style={font=\scriptsize,at={(.02,.98)},anchor=north west},
legend cell align=left
]
\addplot+[black,mark=*,solid,thick]
table[col sep=comma,x=TANNEAL_C,y=DIBL_ELR_ESR_app_mV_perV]
{data/part3_elr_esr_terminal_summary_20261027.csv};
\addlegendentry{Apparent ELR/ESR shift}
\addplot+[black,mark=square*,dashed,thick]
table[col sep=comma,x=TANNEAL_C,y=DrainShift_mV_per_V]
{data/part3_common_current_switching.csv};
\addlegendentry{Measured high-current $V_{\mathrm{SW}}$ shift}
\end{axis}
\end{tikzpicture}
\caption{Two \emph{method-dependent}, non-equivalent
drain-bias shift definitions from the same temperature series.
The solid curve is a mixed ELR/ESR threshold-model extrapolation;
the dashed curve uses the common measured
$|I_D|=3.5\times10^{-4}$ switching current.
Neither curve is the unavailable $10^{-7}$-criterion DIBL.}
\label{fig:part3_secondary_drain_shifts}
\end{figure}

\paragraph{Cross-checks and interpretation.}
The extracted high-drain ON/OFF currents reproduce the
Part~3 metric-table entries at all five temperatures,
providing an independent check that the node pairs
and bias settings were matched correctly.
All reported five-point $R^2$ values are at least
0.99946, and the 3--9-point window ranges are narrow.
The apparent ELR/ESR shift rises from 349.2 to
972.5~mV/V across the tested anneal conditions;
the increase is particularly large at 1000$^\circ$C.

The \emph{other} measured high-current switching
marker tells a more qualified story: it changes from
389.0 to approximately 275.1~mV/V between
800$^\circ$C and 950$^\circ$C, rather than rising
monotonically, before increasing to 895.1~mV/V
at 1000$^\circ$C. Thus, the trend \emph{depends
on the selected extraction definition} over the
800--950$^\circ$C cases, despite both indicators
showing a particularly strong high-drain sensitivity
for the 1000$^\circ$C structure.
In that case the ESR intercept of $-0.9417$~V
is outside the \emph{original} measured
$-0.5$ to $1.2$~V gate interval; it is a model
extrapolate, not an observed current threshold.
Moreover, the separately extended $-1.2$ to
$1.2$~V pilot still did not reach the
original $10^{-7}$ criterion.

The underlying ten \texttt{.plt} source curves and
full consolidated CSV are preserved in the university
Sentaurus workspace, not in this GitHub archive.
Only the precision printed in the provided terminal
output has been transcribed to
\path{data/part3_elr_esr_terminal_summary_20261027.csv}.
Reproducibility scripts
\path{tools/sentaurus_plt_to_idvg.py}
and \path{tools/gm_dibl_exploratory.py}
are included in the repository, but an
independent raw-curve audit requires the ten
original university-server \texttt{.plt} files.
Therefore the ELR/ESR results are
\emph{supplementary numerical estimates}, not
validated subthreshold barrier-lowering magnitudes,
and the strict-current DIBL entries remain N/A.
